\section{System Architecture}
The target system is a stateless, horizontally-scaled API tier written in Haskell, deployed across $N \approx 24$ application nodes behind a network load balancer. Each node serves on the order of 3{,}000 requests per second at peak. The application tier is supported by:
\begin{itemize}[leftmargin=*]
    \item A Redis cluster (six shards, three replicas) acting as the primary fast-path data store for merchant configuration, gateway routing, feature flags, and short-lived transactional state.
    \item A primary relational database (PostgreSQL) for the authoritative transaction ledger.
    \item A risk-decisioning microservice and a fleet of gateway-integration services.
\end{itemize}

\noindent Figure~\ref{fig:baseline-arch} shows the resulting topology and the cross-node call paths that this paper targets.

\begin{figure}[h]
\centering
\resizebox{0.95\textwidth}{!}{%
\begin{tikzpicture}[
    node distance=0.7cm and 1.0cm,
    edge/.style={-Stealth, thick, gray!70!black},
    lb/.style={rectangle, draw, rounded corners, minimum height=0.8cm, minimum width=2.6cm, fill=orange!15, font=\small},
    api/.style={rectangle, draw, rounded corners, minimum height=0.8cm, minimum width=1.8cm, fill=blue!8, font=\footnotesize, align=center},
    store/.style={cylinder, shape border rotate=90, draw, aspect=0.4, minimum height=1.2cm, minimum width=1.8cm, fill=red!8, font=\footnotesize, align=center},
    svc/.style={rectangle, draw, minimum height=0.8cm, minimum width=2.0cm, fill=green!8, font=\footnotesize, align=center}
]
    \node[lb] (lb) {Load Balancer};
    \node[api, below left=1.0cm and 1.7cm of lb] (a1) {API\\Node 1};
    \node[api, below left=1.0cm and -0.1cm of lb] (a2) {API\\Node 2};
    \node[api, below right=1.0cm and -0.1cm of lb] (a3) {$\cdots$};
    \node[api, below right=1.0cm and 1.7cm of lb] (a4) {API\\Node $N$};

    \node[store, below=2.2cm of a2] (redis) {Redis\\Cluster\\(6$\times$3)};
    \node[store, right=of redis] (pg)    {Postgres\\Ledger};
    \node[svc,   right=of pg]    (risk)  {Risk\\Service};
    \node[svc,   right=of risk]  (gw)    {Gateway\\Integrations};

    \draw[edge] (lb) -- (a1);
    \draw[edge] (lb) -- (a2);
    \draw[edge] (lb) -- (a3);
    \draw[edge] (lb) -- (a4);

    \foreach \n in {a1,a2,a3,a4} {
        \draw[edge, dashed] (\n) -- (redis);
    }
    \draw[edge] (a3) -- (pg);
    \draw[edge] (a4) -- (risk);
    \draw[edge] (a4) -- (gw);

    \node[font=\scriptsize, gray, above right=0.05cm and -1.8cm of redis] {5--15 calls / request};
\end{tikzpicture}%
}
\caption{Baseline cluster topology. Each request entering the load balancer is dispatched to one of $N$ stateless API nodes, which together issue 5--15 cross-node calls per request to the shared Redis cluster (dashed) and to the Postgres ledger, risk service, and gateway integrations (solid).}
\label{fig:baseline-arch}
\end{figure}

\section{Request Lifecycle}
A single payment request executes the following stages: (1) authentication and merchant lookup, (2) configuration and routing fetch, (3) risk evaluation, (4) gateway selection, (5) gateway dispatch, and (6) ledger write. Stages (1)--(4) are dominated by configuration-class lookups; stages (5)--(6) involve external calls and authoritative writes. Our caching work targets stages (1)--(4); our contention layer targets stages (5)--(6) and any inventory-style logic embedded within them.

\section{Baseline Performance Characteristics}
Table~\ref{tab:baseline} reports the baseline performance of a representative node under peak load.

\begin{table}[h]
\centering
\begin{tabular}{lr}
\toprule
\textbf{Metric} & \textbf{Value} \\
\midrule
Requests per second per node      & 3{,}000  \\
Cross-node calls per request      & 5--15    \\
p50 API latency                   & 28\,ms   \\
p95 API latency                   & 98\,ms   \\
p99 API latency                   & 285\,ms  \\
Redis p99 latency                 & 12.5\,ms \\
Cross-node calls/sec (cluster)    & $\sim$68{,}000 \\
\bottomrule
\end{tabular}
\caption{Baseline performance of a representative production node before the introduction of caching.}
\label{tab:baseline}
\end{table}

\section{Data Classification}
Configuration data is globally scoped and slowly mutating; routing tables change minutes to hours apart; feature flags rarely; merchant configuration is updated by the merchant via a control-plane API, propagating within seconds. Transactional state, by contrast, is per-request and rapidly mutating. This separation is the basis of our cache-safety framework (Section~\ref{sec:safety}) and of our decision matrix for contention primitives (Chapter~\ref{ch:contention}).

% =============================================================================
